% ============================================================
% A: Problem, motivation, research gap, "same rhythm"   (Joshua, ~2:00)
% ============================================================
\section{Problem and Motivation}
\presenter{Joshua}

\tikzset{
  abox/.style={draw=UCFBlack, rounded corners=3pt, fill=LightGray, minimum height=0.95cm,
               align=center, font=\small, inner sep=5pt},
  aplain/.style={abox, fill=white},
  ahot/.style={abox, fill=UCFGold!60},
  abad/.style={abox, fill=white, draw=SoftGray, text=SoftGray},
  aarr/.style={-{Stealth[length=1.8mm]}, thick, UCFBlack},
  anote/.style={font=\scriptsize, text=SoftGray, align=center},
  aloop/.style={draw=UCFBlack, fill=white, rounded corners=2pt, minimum width=3.1cm,
                minimum height=0.55cm, font=\footnotesize, anchor=west, inner sep=3pt},
}

% ------------------------------------------------------------
\begin{frame}{Find a drum loop by beatboxing it}
  \slidepoint{Sample libraries are searched with text tags. Tags say the genre, not where the kicks, snares and hats land.}

  \centering
  \begin{adjustbox}{max width=\textwidth}
  \begin{tikzpicture}[node distance=0.8cm]
    \node[aplain, minimum width=3.0cm] (q) {\textbf{``boom -- tss -- pf -- tss''}\\[-1pt]
      {\scriptsize 2 bars, beatboxed into a laptop mic}};
    \node[ahot, right=of q, minimum width=2.6cm] (sys) {Rhythm\\search};
    \node[aloop, right=1.0cm of sys, yshift=0.75cm, fill=UCFGold!25] (r1) {1.\ funk loop, kit A \hfill same rhythm};
    \node[aloop, below=0.12cm of r1.south west, anchor=north west] (r2) {2.\ funk loop, kit B \hfill same rhythm};
    \node[aloop, below=0.12cm of r2.south west, anchor=north west, draw=SoftGray, text=SoftGray] (r3) {3.\ soul loop \hfill close};
    \draw[aarr] (q) -- (sys);
    \draw[aarr] (sys.east) -- ++(0.45,0) |- (r1.west);
    \draw[aarr] (sys.east) -- ++(0.45,0) |- (r2.west);
    \draw[aarr] (sys.east) -- ++(0.45,0) |- (r3.west);
    \node[anote, below=0.25cm of q] {\textbf{Input:} an imperfect beatbox};
    \node[anote, below=0.25cm of r3] {\textbf{Output:} drum loops ranked\\by how well the rhythm matches};
  \end{tikzpicture}
  \end{adjustbox}

  \vspace{0.45cm}

  \begin{takeaway}
    People already share grooves by performing them. We want a search box that takes a beat instead of words:
    beatbox 2 bars, get back loops from the library that play the same rhythm, on any drum kit.
  \end{takeaway}

  \note{
    Producers dig through loop libraries using tags like ``funk'' or ``90 BPM''. Those tags describe the
    genre, not where the hits land. Musicians usually communicate a groove by performing it, by beatboxing
    or tapping. So our input is 2 bars of beatbox from a laptop mic, and the output is a ranked list of
    loops that play the same rhythm, whatever kit they're on.
  }
\end{frame}

% ------------------------------------------------------------
\begin{frame}{Same rhythm, completely different sound}
  \slidepoint{A loop and a beatbox of the same groove line up in time, but not in sound.}

  \slidefig[4.3cm]{timbre_gap.pdf}

  \vspace{0.15cm}
  \begin{columns}[T,onlytextwidth]
    \column{0.32\textwidth}
      \textbf{Timbre gap.} A mouth ``pf'' and a kick drum share timing, not spectrum.
    \column{0.32\textwidth}
      \textbf{Imperfect people.} Amateurs rush, drag, drop hi-hats and can't voice every hit.
    \column{0.32\textwidth}
      \textbf{Every voice differs.} So we test on people the model has never heard.
  \end{columns}

  \note{
    Top: one of our library loops. Bottom: a fake beatbox of the same 2 bars. The white lines are beats.
    The hits land in the same places, but the spectra look nothing alike: that's the timbre gap.
    [Optional demo: play the two clips, from results/listening/beatbox\_speed/audio/07\_drums.wav and
    07\_middle.wav.] On top of that, real people are sloppy, and everyone's kick sounds different.
  }
\end{frame}

% ------------------------------------------------------------
\begin{frame}{Research gap}
  \slidepoint{Query-by-beatboxing has been done by transcribing the beatbox first. Nobody has tested learned rhythm embeddings against that.}

  \centering
  \begin{adjustbox}{max width=\textwidth}
  \begin{tikzpicture}[node distance=0.55cm]
    \node[anchor=east, font=\small\bfseries, align=right] (lab1) at (0,0) {Classic\\{\scriptsize\mdseries Kapur et al.;}\\{\scriptsize\mdseries Nakano et al., 2004}};
    \node[abad, right=0.4cm of lab1] (a1) {Beatbox};
    \node[abad, right=of a1] (a2) {Label every\\sound};
    \node[abad, right=of a2] (a3) {Drum\\grid};
    \node[abad, right=of a3] (a4) {Match against\\library grids};
    \draw[aarr, SoftGray] (a1) -- (a2); \draw[aarr, SoftGray] (a2) -- (a3); \draw[aarr, SoftGray] (a3) -- (a4);
    \node[anote, right=0.25cm of a4, text=red!70!black, align=left] {one misheard\\sound breaks\\the match};

    \node[anchor=east, font=\small\bfseries, align=right] (lab2) at (0,-1.8) {Ours};
    \node[aplain, right=0.4cm of lab2] (b1) {Beatbox};
    \node[ahot, right=of b1, minimum width=3.6cm] (b2) {Learned\\rhythm embedding};
    \node[aplain, right=of b2] (b3) {Nearest loops,\\any kit};
    \draw[aarr] (b1) -- (b2); \draw[aarr] (b2) -- (b3);
    \node[anote, below=0.1cm of b2] {trained so beatbox and loop of\\the same rhythm land together};
  \end{tikzpicture}
  \end{adjustbox}

  \vspace{0.15cm}
  \raggedright\footnotesize
  \begin{itemize}
    \item \textbf{Close but different:} TRIA (2025) generates drums from beatboxes (no search); Ramires et al.\ (2018)
          adapt \emph{transcription} to each user; CLAP captures what audio means, not where hits land.
    \item \textbf{Our question:} do learned rhythm embeddings beat transcribe-then-match on people and drum kits never seen in training?
  \end{itemize}

  \note{
    The idea goes back to 2004: Kapur and colleagues, and Nakano and colleagues, transcribed each beatbox
    sound into a drum hit, then matched symbolically. If one sound is misheard, the match breaks. Recent work
    like TRIA uses beatboxes to generate drums rather than search, and per-user adaptation has only been
    done for transcription. So our gap: compare a learned rhythm embedding against transcribe-then-match,
    on held-out people and held-out kits.
  }
\end{frame}

% ------------------------------------------------------------
\begin{frame}{What counts as ``the same rhythm''?}
  \slidepoint{Each loop becomes a 3-voice grid. Same grid = same rhythm.}

  \slidefig[2.6cm]{rhythm_grid.pdf}

  \vspace{0.1cm}
  \begin{columns}[T,onlytextwidth]
    \column{0.55\textwidth}
      \small
      \begin{itemize}
        \item \textbf{3 voices:} kick, snare, hats (cymbals count as hats). Toms, crashes and ghost notes are left out.
        \item \textbf{12 steps per beat} hold both straight 16ths and triplets, so a swung beat is a different rhythm from a straight one.
      \end{itemize}
    \column{0.40\textwidth}
      \begin{takeaway}
        \centering
        {\Large\bfseries 8,626}\\
        distinct rhythms in 8,643 two-bar loops\\[2pt]
        {\footnotesize\color{SoftGray}Almost every loop is its own pattern, so each query has only a few right answers.}
      \end{takeaway}
  \end{columns}

  \note{
    This is the grid for the loop on the last slide. As the music lead, I wrote these rules: three voices,
    and twelve steps per beat, which is fine enough to tell straight 16ths from triplet swing. Two clips
    are the same rhythm only if the grids match exactly. Because both sides come from the same MIDI,
    the answer key is free. And 8,626 of the 8,643 loops are their own pattern, so this is close to
    finding one exact loop in the library.
  }
\end{frame}
